\documentclass[letterpaper,11pt]{article}

\usepackage[T1]{fontenc}
\usepackage[utf8]{inputenc}
\usepackage[bitstream-charter]{mathdesign}
\usepackage[top=0.7in,bottom=0.7in,left=0.9in,right=0.9in]{geometry}
\usepackage[hidelinks]{hyperref}
\usepackage{parskip}
\linespread{1.04}
\setlength{\parskip}{7pt}
\pagestyle{empty}
\raggedright

\begin{document}

%----------HEADER----------
\begin{center}
    {\fontsize{18}{20}\selectfont\bfseries DEVA ANAND}\\ \vspace{3pt}
    {\small
    \href{mailto:devaanand@umass.edu}{devaanand@umass.edu}
    \enspace$\vert$\enspace
    (413) 315-1715
    \enspace$\vert$\enspace
    \href{https://devaanand.com}{devaanand.com}
    \enspace$\vert$\enspace
    \href{https://github.com/Deva-1903}{github.com/Deva-1903}}
\end{center}
\vspace{6pt}

\noindent July 15, 2026

\vspace{2pt}

\noindent Cloudflare\\
Austin, TX

\vspace{2pt}

\noindent \textbf{Re: Software Engineering Intern}

\vspace{4pt}

\noindent Dear Cloudflare Hiring Team,

I have wanted to work at Cloudflare since the first time I deployed to Workers and watched a full application run at the edge without provisioning a single server. I am a Computer Science master's student at UMass Amherst applying for the Software Engineering Intern role, and Cloudflare is one of the few companies where the infrastructure I already build on for fun is also a mission I believe in: a faster, more secure, and genuinely open Internet.

That is not abstract admiration. My project cf-ai-research-scout is built end to end on your developer platform: an AI research agent that runs on Cloudflare Workers, holds stateful conversations through Durable Objects, persists to D1, and coordinates multi-step research with WorkflowEntrypoint so tasks survive process restarts and individual tool failures without any external database. Building it convinced me that durable, edge-native workflows are the right way to make AI systems reliable, and it is the closest I have come to developing at Internet scale on my own.

I am also bringing Cloudflare deeper into my current work. As the sole engineer on a research platform for UMass's Culture and Morality Lab, I am planning to move our file and artifact storage onto Cloudflare R2, both for its S3-compatible ergonomics and, honestly, because zero egress fees are a real win on an academic budget. It is a small decision, but it is exactly the "spot the friction, reach for the better tool" instinct your team describes, and it is a pattern I keep repeating: pick the platform that lets me ship the reliable thing fastest.

The rest of my background lines up with how you build. At Wysa I owned backend security and reliability for a multi-tenant service used by 8+ UK clinical clients, implementing encryption for sensitive data and remediating 20+ security findings, because secure and reliable was the product. In my own projects I am AI-native by default: I hardened an agentic search service with 150+ automated tests and grounded its LLM output in verifiable, cell-level provenance rather than trusting a black box. I like getting things done, I like shipping, and I like sharing what I learn along the way.

I would be thrilled to spend an internship pushing code that reaches hundreds of millions of people and learning from the engineers who keep that network fast and honest. Thank you for considering my application, and I would welcome the chance to talk.

\vspace{6pt}

\noindent Sincerely,\\
Deva Anand

\end{document}
